\section{Two Incompatible and Irreconcilable Worldviews}

There are two incompatible and irreconcilable worldviews.

\paragraph{Anti-tradition}
One, is the idea that humans evolved from animals and then, even within human society, there is a continual process of evolution and progress to higher stages of intellectuality and morality.

\paragraph{Tradition}
The second is that humans began in a perfected state, or ``Golden Age", but through choices there has been a continual state of decline from that state of perfection, and only a cataclysmic event can restore a new state of perfection.

\paragraph{Synthesis}
There is certainly no scientific basis for evolution or progress. If, in fact, we do seem to observe change, progress, or evolution, then its etiology lies elsewhere. For every idea or movement, there is an impetus that sets it in motion, that brings about its realization. However, over time, that initial impetus gets distorted and every movement or organisation will tend to decline even to the point of turning into its very opposite. That is due to the common unconscious state of humanity. Nevertheless, it is possible at certain key moments in the process — and time is qualitative, not quantitative == that another conscious impetus can be applied to bring the process back into alignment. This can only come about through the activities of a conscious elite.



\flrightit{Posted on 2010-09-12 by Cologero }

\begin{center}* * *\end{center}

\begin{footnotesize}\begin{sffamily}



\texttt{James O'Meara on 2010-09-13 at 11:53 said: }

Guenon, if I understand him, thought that the movement within each cycle was the product of a confluence of opposite currents, ascending and descending, as it were. But I don't think the positive movements were able to completely reverse or halt the overall cyclical movement. 

At the end of the cycle [Judgment Day] the positive remainder [the Qabalistic klippoth] would be passed on to the new cycle; this was the function of the elite, for Evola. But I don't believe even he thought there could be an ultimate, or permanent, ``reversal" of the cycle. At best, the elite might survive the cataclysm and pass into the new cycle [the Earthly Paradise].


\hfill

\texttt{Cologero on 2010-09-13 at 19:36 said: }

``That is due to the common unconscious state of humanity"

You could start by Evola's description in The Individual and the Becoming of the World. There are several references on Gornahoor.

Or Louis-Claude de Saint-Martin on the Man of Torrent (In the stream): Men without will-power and little individuality, who follow the way of life in a given epoch without bothering about anything except their elementary needs.

Gurdjieff on the man as asleep.

Guenon on men dominated by a modern ideology who are unable to recognize higher dimensions.

You take on the moniker of a Hermetist: How do you view the Great Work?


\hfill

\texttt{Cologero on 2010-09-13 at 19:57 said: }

There are movements that have their own course within a part of the entire cycle. Guenon's model gets complex as he described in the Symbolism of the Cross, but there the point is not the regeneration of the Earth, but an entirely new world. I'm not sure what you mean by currents. Guenon has written about the gunas, but there are three, not two.

Nevertheless, there are still possibilities within a cycle\footnote{I addressed Guenon's viewpoint here: \url{http://www.gornahoor.net/?p=35}.}.


\hfill

\texttt{James O'Meara on 2010-09-14 at 18:34 said: }

``I'm not sure what you mean by currents. Guenon has written about the gunas, but there are three, not two."

I can't find the quote right now, but I think he was in fact talking about the gunas as features of cosmology. However, even with three or more gunas, they would act together in various combinations or `strengths’ [today, for example, raja would hardly dominate] to produce an overall direction or vector. If the cyclical descent could be compared to a curling stone, the various forces, gunas, magical chains such as the UR group proposed, counter-traditional forces, etc, would be the sweepers in front, speeding up or retarding, perhaps even slightly reversing the movement of the stone, but never bringing it to a halt or sending it back up.


\end{sffamily}\end{footnotesize}
